\section{Conclusion}
\label{sec:conclusion}

This work introduced \textsc{SenseMath}, a controlled benchmark that measures prompt-sensitive shortcut exploitation on matched item families.
By pairing strong-shortcut, weak-shortcut, and control variants within each family, the benchmark isolates the prompting asymmetry from confounds such as general prompt quality or item difficulty.

Our experiments across five models yield three main findings.
First, a category-agnostic NS prompt produces an asymmetric effect at $d{=}4$: accuracy on shortcut-amenable items is maintained or improved while accuracy on matched controls drops, most consistently for GPT-4o-mini and Qwen3-8B.
Second, per-category radar profiles reveal that models show selective strengths and reversals across shortcut types. Qwen models lead on magnitude estimation and structural shortcuts, while GPT-4o-mini dominates cancellation, equation reasoning, and option elimination. This establishes number sense as a multi-dimensional construct rather than a single ability.
Third, a three-level evaluation (Use, Judge, Generate) reveals crossed dissociations: models that exploit shortcuts well do not necessarily recognise when shortcuts apply, and vice versa.

\paragraph{Limitations.}
Several limitations qualify these findings.
We have not conducted a human behavioural study to anchor LLM performance against human number sense; such a study is planned but not yet complete.
The benchmark contains 50 families per category, which is sufficient for detecting moderate-to-large effects but limits statistical power for fine-grained within-category analyses.
The Strict prompting condition produces truncated responses for open-weight models at 512 tokens, restricting its use as a negative control to GPT-4o-mini.
The multiple-choice format constrains the space of possible responses; free-response evaluation would provide richer signal about the reasoning process.
Finally, the 8 categories, while grounded in the mathematics education literature, are not exhaustive. Other shortcut types (e.g., modular arithmetic, dimensional analysis) remain untested.

\paragraph{Future directions.}
Several extensions follow naturally from this work.
A human validation study would establish whether the shortcut profiles we observe in LLMs mirror human number-sense behaviour.
Moving from multiple-choice to free-response items would allow finer-grained analysis of how models construct shortcut strategies.
Evaluating additional models, particularly reasoning-specialised systems and larger open-weight models, would clarify how number sense scales with model capability and training data.
Finally, the asymmetric prompting effects documented here suggest that shortcut-sensitive reasoning is in principle trainable; targeted fine-tuning on curated shortcut solutions could shift models from algorithmic defaults toward more flexible numerical reasoning.
